\section{Opportunities and Limitations of AI Agents in Software Development}

Across both parts, the project shows that AI coding agents can do substantial, genuinely useful work. They produced complete features spanning backend and frontend, generated systematic test coverage for non-trivial scenarios, and --- once given a fixed, narrow checklist instead of an open-ended instruction to "review the code" --- reliably found problems a human reviewing informally had missed, such as untested code paths or inconsistent naming. This points to a real opportunity: an agent cannot replace the judgment a human reviewer brings, but a large share of the mechanical, checklist-style part of quality review can be handed to an agent built specifically for that job, leaving human attention for the judgment calls a checklist cannot make on its own.

Structure did not remove the need for human judgment in Part 2; it relocated it. Where Part 1 required judgment scattered throughout implementation, Part 2 concentrated it at fewer points, primarily plan approval --- a shift in when and how judgment is used, not its elimination. This success is also partly self-referential: several of the Charter's rules were written directly from mistakes already observed in Part 1, so how well the structured approach would generalize to a project without that firsthand knowledge remains untested.

\textbf{Order/learning effect.} A further limitation concerns the order in which the two parts were carried out. Part 1 was built first, and Part 2 began only once it was finished, by the same person. By the time Part 2 started, general familiarity with prompting these agents, with the application's own domain, and with where earlier mistakes had come from had already accumulated --- none of which existed when Part 1 began. This means part of Part 2's improvement cannot be cleanly separated from ordinary learning over the course of the project, rather than attributed to the structural changes alone. Separating the two effects would require a different study design --- for example, running the structured approach first and the unstructured approach second, or repeating the comparison with a different person. This thesis does neither: it runs each approach once, in a fixed order, so the two effects cannot be told apart here.

\textbf{Tool and environment as a second variable.} The comparison in this thesis deliberately holds the underlying model constant, Claude Sonnet 4.6 throughout, specifically so that differences between Part 1 and Part 2 could be attributed to collaboration structure rather than model capability. The development environment, however, was not held constant alongside it: Part 1 was built using Claude Code inside Visual Studio Code, while Part 2 moved to Google Antigravity. This means the comparison carries a second variable beyond structure --- the tool itself --- riding alongside the one the thesis set out to isolate. Whether Part 2's outcomes would hold if the same structured methodology had instead been run inside Part 1's original environment, without also changing tools, is a question this project's design cannot answer on its own.

